\documentclass[11pt,a4paper]{article}
\usepackage[brazil]{babel}
\usepackage[utf8]{inputenc}
\usepackage[T1]{fontenc}
\usepackage{lmodern}
\renewcommand{\baselinestretch}{1.05}
\usepackage{mhchem}
\usepackage{graphicx}
\usepackage{svg}
\usepackage{amssymb}
\usepackage{amsmath}
\usepackage{enumitem}
\usepackage{geometry}
\usepackage{booktabs}
\geometry{top=2cm, bottom=2cm, left=2.5cm, right=2.5cm}

\newcommand{\oC}{$^\circ$C}
\newcommand{\e}[1]{$\times 10^{#1}$}
\newcommand{\moll}{~mol~L$^{-1}$}

\setlist[itemize,1]{leftmargin=*, itemsep=2pt, topsep=2pt}
\setlist[enumerate,1]{leftmargin=*, itemsep=2pt, topsep=2pt}

\pagestyle{empty}

\begin{document}

\begin{center}
  {\LARGE\bfseries QG101 -- Química Geral}\\[4pt]
  {\large\bfseries Bloco 9 -- Forças Intermoleculares e Estados da Matéria}\\[2pt]
  \textit{Material didático -- Aula Teórica 9}
\end{center}

\bigskip

% ================================================================
\section*{1. Ligações intramoleculares x interações intermoleculares}

As ligações \textbf{intramoleculares} (iônicas e covalentes, vistas nos blocos
anteriores) mantêm os átomos unidos \emph{dentro} de uma molécula ou de um
sólido iônico. São fortes: tipicamente entre 150 e 1000~kJ/mol.

As \textbf{interações intermoleculares} (ou forças de van der Waals, em
sentido amplo) atuam \emph{entre} moléculas vizinhas. São muito mais fracas
-- de poucos kJ/mol a, no máximo, cerca de 40~kJ/mol -- mas são responsáveis
por propriedades macroscópicas essenciais: ponto de fusão, ponto de
ebulição, viscosidade, tensão superficial e o próprio estado físico
(sólido, líquido ou gás) de uma substância em condições ambientes.

\begin{quote}
\textit{Quando a água entra em ebulição, as ligações covalentes O--H
\textbf{não} são rompidas -- apenas as ligações de hidrogênio
\textbf{entre} as moléculas de \ce{H2O} são vencidas.}
\end{quote}

% ================================================================
\section*{2. Tipos de forças intermoleculares}

\subsection*{2.1 -- Forças de London (dispersão)}

Presentes em \textbf{todas} as moléculas e átomos, mesmo os apolares. Surgem
porque a nuvem eletrônica flutua continuamente, criando dipolos
\textbf{instantâneos} que induzem dipolos correspondentes nas moléculas
vizinhas (Figura~1).

A intensidade das forças de London cresce com a \textbf{polarizabilidade}
da nuvem eletrônica, que por sua vez aumenta com o número de elétrons
(tamanho/massa molar) e com a área de contato entre as moléculas (moléculas
alongadas interagem mais que moléculas compactas de mesma massa molar).

\begin{center}
\begin{tabular}{lccc}
\toprule
Gás nobre & \ce{He} & \ce{Ne} & \ce{Ar} \\
\midrule
Ponto de ebulição (\oC) & $-269$ & $-246$ & $-186$ \\
\bottomrule
\end{tabular}
\end{center}

O aumento do ponto de ebulição de \ce{He} a \ce{Ar} ilustra o efeito da
polarizabilidade: mais elétrons $\Rightarrow$ nuvem mais deformável
$\Rightarrow$ forças de London mais intensas. O mesmo padrão explica por
que, à temperatura ambiente, \ce{F2} e \ce{Cl2} são gases, \ce{Br2} é
líquido e \ce{I2} é sólido.

\subsection*{2.2 -- Forças dipolo-dipolo}

Ocorrem entre moléculas que possuem \textbf{dipolo permanente}
(momento de dipolo $\mu \neq 0$): a extremidade $\delta^+$ de uma molécula
atrai a extremidade $\delta^-$ da vizinha (Figura~1). São, em geral, mais
intensas que as forças de London entre moléculas de massa molar
semelhante.

\textbf{Exemplo:} o butano (\ce{C4H10}, $M = 58$~g/mol, apolar) ferve a
$-1$\oC, enquanto a acetona (\ce{C3H6O}, $M = 58$~g/mol, polar) ferve a
$56$\oC. Mesma massa molar, mas a acetona soma forças dipolo-dipolo às
forças de London, exigindo mais energia para separar as moléculas.

\subsection*{2.3 -- Ligações de hidrogênio}

Caso particular e especialmente forte de atração dipolo-dipolo, que ocorre
quando um átomo de \textbf{H} está covalentemente ligado a um átomo pequeno
e muito eletronegativo (\ce{N}, \ce{O} ou \ce{F}) e interage com um par de
elétrons não-ligantes de um \ce{N}, \ce{O} ou \ce{F} \emph{vizinho}
(Figura~1). Energia típica: 10--40~kJ/mol -- mais forte que uma
dipolo-dipolo comum, porém ainda muito mais fraca que uma ligação covalente.

\begin{center}
\begin{tabular}{lcccc}
\toprule
Hidreto (grupo 16) & \ce{H2O} & \ce{H2S} & \ce{H2Se} & \ce{H2Te} \\
\midrule
Ponto de ebulição (\oC) & $100$ & $-60$ & $-41$ & $-2$ \\
\bottomrule
\end{tabular}
\end{center}

Do \ce{H2S} ao \ce{H2Te} o ponto de ebulição \textbf{sobe} com o aumento de
tamanho/polarizabilidade (forças de London), como esperado. Mas a
\ce{H2O} é a grande anomalia da série: \textbf{sem} ligações de hidrogênio,
seu ponto de ebulição seria esperado em torno de $-90$\oC; com elas, sobe
para $100$\oC.

% ================================================================
\section*{3. Resumo comparativo}

\begin{center}
\begin{tabular}{lp{5.4cm}cl}
\toprule
Força & Origem & Energia (kJ/mol) & Exemplo \\
\midrule
London (dispersão) & Dipolos instantâneos induzidos & $0,05$--$40$ &
  \ce{Ar}, \ce{CH4}, \ce{I2} \\
Dipolo-dipolo & Atração entre dipolos permanentes & $5$--$25$ &
  Acetona, \ce{SO2} \\
Ligação de hidrogênio & \ce{H} (em N--H, O--H, F--H) $\cdots$ par livre &
  $10$--$40$ & \ce{H2O}, \ce{NH3}, \ce{HF} \\
\midrule
\textit{(referência)} Ligação covalente & Compartilhamento de elétrons &
  $150$--$1000$ & \ce{O-H}, \ce{C-C} \\
\bottomrule
\end{tabular}
\end{center}

\bigskip
\begin{center}
\includesvg[width=0.95\textwidth]{figuras/bloco09_fig1_forcas}

{\small Figura 1 -- Representação esquemática das três principais forças
intermoleculares.}
\end{center}

\clearpage

% ================================================================
\section*{4. Influência nas propriedades físicas}

Quanto mais intensas as forças intermoleculares, \textbf{mais energia} é
necessária para separar as moléculas -- logo, maiores tendem a ser:

\begin{itemize}
\item o \textbf{ponto de fusão} e o \textbf{ponto de ebulição};
\item a \textbf{viscosidade} (resistência ao escoamento: glicerina e mel,
  com várias hidroxilas por molécula, formam extensas redes de ligações de
  hidrogênio e escoam muito mais lentamente que a água ou a gasolina);
\item a \textbf{tensão superficial} (a água, por suas ligações de
  hidrogênio, tem tensão superficial bem maior que solventes orgânicos
  comuns, permitindo, por exemplo, que insetos caminhem sobre sua
  superfície).
\end{itemize}

\begin{center}
\begin{tabular}{lcccc}
\toprule
Substância & \ce{Ne} & \ce{CH4} & \ce{NH3} & \ce{H2O} \\
Massa molar (g/mol) & $20$ & $16$ & $17$ & $18$ \\
Força predominante & London & London & dipolo-dipolo / H & ligação de H \\
Ponto de ebulição (\oC) & $-246$ & $-161$ & $-33$ & $100$ \\
\bottomrule
\end{tabular}
\end{center}

Note que \ce{CH4} e \ce{H2O} têm massas molares muito próximas, mas pontos
de ebulição radicalmente diferentes: o que importa não é apenas o
\emph{tamanho} da molécula, mas o \textbf{tipo} de força intermolecular
que ela pode formar.

% ================================================================
\section*{5. Estados físicos e mudanças de estado}

Nos sólidos, as forças intermoleculares mantêm as partículas em posições
fixas (vibrando no lugar); nos líquidos, as partículas deslizam umas sobre
as outras mas permanecem próximas; nos gases, as forças intermoleculares
são vencidas pela energia cinética e as partículas se movem livremente,
muito afastadas entre si.

Ao fornecer calor a um sólido (Figura~2), a temperatura sobe até atingir o
\textbf{ponto de fusão}, onde permanece \textbf{constante} enquanto o
sólido vira líquido: toda a energia fornecida é usada para vencer parte
das forças intermoleculares da rede, não para aumentar a energia cinética
média (e, portanto, a temperatura). O mesmo ocorre no \textbf{ponto de
ebulição}. Por isso a \emph{temperatura não é um bom indicador} de quanto
calor já foi fornecido durante uma transição de fase.

\begin{center}
\includesvg[width=0.62\textwidth]{figuras/bloco09_fig2_curva_aquecimento}

{\small Figura 2 -- Curva de aquecimento típica da água sob pressão
constante: a temperatura permanece constante durante as mudanças de fase.}
\end{center}

% ================================================================
\section*{6. Propriedades anômalas da água}

A rede tridimensional de ligações de hidrogênio (cada molécula de
\ce{H2O} pode doar 2 e aceitar 2 ligações de hidrogênio) explica diversas
propriedades incomuns da água:

\begin{itemize}
\item \textbf{Ponto de ebulição elevado} para uma molécula tão pequena
  (Seção~2.3).
\item \textbf{O gelo é menos denso que a água líquida}: no gelo, as
  ligações de hidrogênio organizam as moléculas em uma estrutura
  hexagonal aberta, com mais espaço vazio do que no líquido (onde a rede
  é parcial e mais compacta). Por isso o gelo \textbf{flutua} -- um
  comportamento raro entre as substâncias, e essencial para a vida
  aquática em climas frios (a camada de gelo isola a água líquida abaixo).
\item \textbf{Calor específico elevado} ($4{,}18$~J~g$^{-1}$\oC$^{-1}$):
  romper parcialmente a rede de ligações de hidrogênio para aumentar a
  energia cinética média consome bastante energia, fazendo da água um
  excelente regulador térmico (moderação climática costeira, regulação da
  temperatura corporal).
\item \textbf{Tensão superficial elevada}, já mencionada na Seção~4.
\end{itemize}

% ================================================================
\section*{7. Síntese}

\begin{center}
\begin{tabular}{p{3.4cm}p{5.2cm}p{5.6cm}}
\toprule
Conceito & Ideia central & Consequência prática \\
\midrule
Intra x intermolecular & Ligações fortes \textit{dentro}; forças fracas
  \textit{entre} moléculas & Fusão/ebulição rompem apenas interações
  intermoleculares \\
London & Dipolos instantâneos; presentes sempre & Cresce com massa
  molar/polarizabilidade \\
Dipolo-dipolo & Dipolos permanentes & Soma-se a London em moléculas
  polares \\
Ligação de H & H ligado a N/O/F $\cdots$ par livre & Pontos de
  ebulição anomalamente altos (\ce{H2O}, \ce{NH3}, \ce{HF}) \\
Estados da matéria & Energia cinética x forças intermoleculares & Define
  sólido/líquido/gás e as transições entre eles \\
\bottomrule
\end{tabular}
\end{center}

\bigskip
\noindent\textit{A lista de exercícios correspondente a este bloco está no
arquivo \texttt{bloco09\_lista\_exercicios.pdf}.}

\end{document}
